\chapter{1977年江苏镇江地区卷}

\begin{problem}
计算： （1）\((-1)^{4}-8\times(-1\frac{1}{2})^{3}+20.1\div(-0.1^{2})\)；（2）\(\frac{1}{2\sqrt{2}} - \frac{2\sqrt{3}}{3 - \sqrt{3}} - \frac{1}{2} (\sqrt{\frac{1}{2}} - \sqrt{12})\)；（3）\( \lg 1 - \lg 0.001 + \log_{2}(-2)^{4} + \log_{3}\sqrt{3} \)；（4）当 \( a = \sqrt{10.1}, b = 8 \) 时，求式子 \( \frac{(a^{2}b^{-1})^{8}\sqrt{a b^{\frac{2}{3}}}}{a^{2}b^{-4}} \) 的值. 

\end{problem}

\begin{problem}
（1）求 \( \sin75^{\circ}+\frac{1}{2}\cos(-585^{\circ})\cdot tg240^{\circ} \) 的值. （2）证明 \( \frac{2}{\operatorname{tg} x + \operatorname{ctg} x} = (\sin x + \cos x)^{2} - 1 \) .

\end{problem}

\begin{problem}
解方程 \(\frac{x - 1}{x^2 + 4x + 3} - \frac{2}{x^2 - 1} = 0\). 

\end{problem}

\begin{problem}
求下列各函数的定义域（自变量 \(x\) 的取值范围）： （1）\( y = -x + 3 \)；（2）\( y = \frac{1}{x^2 + 2x - 3} \)，（3）\( y = \sqrt{-2x + 10} \)；（4）\( y = \lg (-x^2 + 5x - 4) \).

\end{problem}

\begin{problem}
从圆外一点引圆的两条割线，试证一条割线和它的圆外部分的乘积等于另一条割线和它圆外部分的乘积. 

\end{problem}

\begin{problem}
为了测量底部不能直接到达的高塔 \(AB\) 的高，在地面上引一条基线 \(CD=60\mathrm{m}\)（\(B\)、\(C\)、\(D\) 不在一直线上），今测得 \( \angle ACB=30^{\circ} \) ，

\( \angle BCD=75^{\circ},\angle BDC=45^{\circ} \) ，求 \(AB\). 

\end{problem}

\begin{problem}
（1）画出抛物线 \( y = x^{2} + 2 \) 的图象，并写出它的顶点坐标、对称轴的方程和极值；（2）如果上述抛物线的顶点作为新的坐标原点 \( O_{1} \) ，把它的对称轴作为新的纵轴 \( O_{1}Y_{1} \) ，求这抛物线在新的直角坐标系 \( X_{1}O_{1}Y_{1} \) 中的方程，并写出它的准线方程和焦点坐标. 

\end{problem}

\begin{problem}
某工厂四年来的产量，第一年到第三年的产量是按等差数列增长的，三年的总产量共1500吨；第二年到第四年

的产量是按等比数列增长的，这三年的产量共1820吨. 求每年的产量各是多少吨？

\end{problem}

\begin{problem}
三角形 \(ABC\) 中，\(AE\) 和 \(BD\) 相交于 \(G\)（如图），已知 \(AD:DC=3:1\)，\(BG=GD\)，求：\( \frac{AG}{GE} \) .
\end{problem}

